\documentclass[10pt,a4paper]{amsart}
\usepackage[margin=19mm]{geometry}
\usepackage{amsmath,amssymb,mathtools}
\usepackage[scr=rsfs,cal=euler]{mathalfa}
\usepackage{xfrac}
\usepackage[unicode,hidelinks,bookmarks=false]{hyperref}
\newcommand{\ie}{i.e.\ }
\newcommand{\cf}{cf.\ }
\newcommand{\resp}{respectively\ }
\newcommand{\Subsection}{\subsection}
\providecommand{\nobreakdash}{\nobreak\hbox{-}}
\setlength{\emergencystretch}{2em}
\multlinegap0pt
\usepackage{amssymb}
\usepackage{tikz}
\usepackage{mathtools}
\newtheorem{lemma}{Lemma}
\def\E{\mathbb E}
\title{Oscillations of random multiplicative functions under initial bias}
\author{Rodrigo Angelo, Max Wenqiang Xu}
\date{Source: arXiv:2411.14447v3 (CC BY 4.0)}
\begin{document}
\maketitle

\noindent\emph{Source and licence.} Oscillations of random multiplicative functions under initial bias. Version arXiv:2411.14447v3 (CC BY 4.0), distributed by arXiv under the Creative Commons Attribution 4.0 International licence, http://creativecommons.org/licenses/by/4.0/, as recorded in the arXiv licence metadata for this exact version and shown on the abstract page https://arxiv.org/abs/2411.14447v3. Full licence notice: \texttt{licenses/arxiv-2411.14447-CC-BY-4.0.txt}.\par\smallskip

\noindent\textit{Editorial scope.} Counted unit: the article's lemma transition (source0.tex lines 218--221). The article's complete original proof of it is delivered with it (source0.tex lines 223--260, one \texttt{proof} environment of 38 lines). No mathematics is written, repaired or re-derived: the statement and the proof are the article's own lines, in the article's own notation, and no retained line is substituted or re-flowed.  No further source-local context is needed: every cross-reference of every retained line resolves inside the two delivered spans.  Nothing was omitted from any retained span, so the package carries no omission record.  No retained line contains a citation, so the wrapper carries no bibliography and no bibliography entry.  No retained line contains a figure environment, an image-inclusion command or drawing code, so no image is delivered and the package declares no asset dependency.  Editorial formatting only, in addition: an AMS \texttt{amsart} preamble that restates the article's own declarations and macros used by the retained lines, namely the environments \texttt{lemma} on separate counters, together with the standard packages those lines need.  The retained environments print in this document as Lemma 1 in order of appearance, whereas the article prints the same items with the numbering of its own full document.  The complete native source of the article as distributed in the arXiv e-print for this exact version is bundled unchanged as \texttt{sources/168913/source0.tex}.
\begin{lemma}\label{transition} 
     Let $f$ denote the random completely multiplicative function, conditioned on $f(p) = 1$ for $p \le y$. Assume $\frac{20\log \log x}{\log x}\le t \le \frac{1}{y^{\frac{1}{2}}}$. Then for $s = \frac{1}{2} + t$, with probability $1-o(1)$ one has
$$\sum_{n \le x} \frac{f(n)}{n^s} = \exp(R(t)+O(1)).$$
\end{lemma}

\begin{proof}
We begin by showing
$$
\mathbb{E} \left| \sum_{n \leq x} \frac{f(n)}{n^s}  \prod_p \left( 1 - \frac{f(p)}{p^s} \right) -1\right|^2 =o(1).
$$
By writing the finite sum as the entire sum minus its tail, and using the Euler product we get that this is equal to
$$\mathbb{E} \left| \sum_{n > x} \frac{f(n)}{n^s}  \prod_{p} \left( 1 - \frac{f(p)}{p^s} \right) \right|^2.$$ 
    Factoring out $p \le y$ gives that
\begin{equation}\label{eqn: full}
    \prod_{p \le y} \left(1  - \frac{1}{p^s}\right)^2\mathbb{E} \left| \sum_{n > x} \frac{f(n)}{n^s}  \prod_{p>y} \left( 1 - \frac{f(p)}{p^s} \right) \right|^2.    
    \end{equation}
Now notice that $\E[f(mn)]=1$ precisely when $mn$ is of the shape $ab^2$ where all the prime factors of $a$ are $\le y$ and the prime factors of $b$ are $>y$. Otherwise $\E[f(mn)] = 0$.
    
    Expand the inner sum and consider the coefficient of $f(N)$ for each $N = ab^2$ as described. We can write $ab^2=n_1n_2m_1m_2$, where $n_1,n_2$ are from the sum $\sum_{n>x}\frac{f(n)}{n^{s}}$ and $m_1,m_2$ from the infinite product. Since the prime factors of $a$ are all $\le y$, we may further write $n_1=d_1\ell_1$ and $n_2=d_2\ell_2$ where $d_1d_2=a$, and $m_1m_2\ell_1\ell_2=b^2$. As such, the number of times $f(N)$ with $N=ab^2$ appears in the expansion is at most $\tau_2(a)\tau_4(b^2)$.
Together with an application of Rankin's trick and $s = \frac{1}{2}+t$, this implies that  $\mathbb{E} \left| \sum_{n > x} \frac{f(n)}{n^s}  \prod_{p>y} \left( 1 - \frac{f(p)}{p^s} \right) \right|^2$
\begin{equation*}
    \begin{split}
    \le \sum_{\substack{ab^2>x^2 \\ p|a \implies p\le y\\ p|b \implies p>y}} \frac{\tau_2(a)\tau_4(b^2)}{(ab^2)^s} &\ll \frac{1}{x^t}\sum_{\substack{p|a \implies p\le y\\ p|b \implies p>y}} \frac{\tau_2(a)\tau_4(b^2)}{(ab^2)^{\frac{1}{2}+\frac{t}{2}}} \\ & \ll \frac{1}{x^t}\prod_{p \le y} \left(1-\frac{1}{p^{\frac{1}{2}+\frac{t}{2}}}\right)^{-2} \prod_{p>y} \left(1-\frac{1}{p^{1+t}}\right)^{-10}.     
    \end{split}
\end{equation*}
    Given this, the contribution of all primes $p \le y$ in \eqref{eqn: full} is 
    $$\prod_{p \le y} \left(1  - \frac{1}{p^{\frac{1}{2}+t}}\right)^2\left(1-\frac{1}{p^{\frac{1}{2}+\frac{t}{2}}}\right)^{-2} \ll \exp \left(t\sum_{p \le y} \frac{\log p}{p^{\frac{1}{2}}}\right) \ll \exp(t y^{\frac{1}{2}}).$$
 The contribution from primes $p>y$ is 
    $$\prod_{p>y} \left(1-\frac{1}{p^{1+t}}\right)^{-10} \le \zeta(1+t)^{10} \ll \frac{1}{t ^{10}}.$$
    Therefore, \eqref{eqn: full} is 
    $$\ll \frac{\exp(ty^{\frac{1}{2}})}{x^{t}t^{10}}=o(1).$$
    Note that $x^{t} \gg (\log x)^{20}$, $t^{10} \ge \frac{1}{(\log x)^{10}}$, and $\exp(ty^{\frac{1}{2}}) \ll 1$,  using $\frac{20\log \log x}{\log x}\le t \le \frac{1}{y^{\frac{1}{2}}}$.

  Given the above second moment estimate, we can see that with probability at least $1 - o(1)$ one has
    $$\frac{1}{2}\prod_{p} \left( 1 - \frac{f(p)}{p^s} \right)^{-1} \le \sum_{n \le x} \frac{f(n)}{n^s} \le 2\prod_{p} \left( 1 - \frac{f(p)}{p^s} \right)^{-1}.$$
    The lemma then follows from the fact that
    $$\prod_{p} \left( 1 - \frac{f(p)}{p^s} \right)^{-1} = \exp(R(t)+O(1)),$$
via the approximation $\left(1-\frac{f(p)}{p^s}\right)^{-1} = \exp\left(\frac{f(p)}{p^s}+\frac{1}{2p^{2s}}+O\left(\frac{1}{p^\frac{3}{2}}\right)\right)$. 

    

    
\end{proof}

\end{document}
